% Folder 77, batch 04 (archivists' pages 61-80).
% Pass: Opus 5.5 (claude-opus-5-5), 2026-10-03 — first pass, unchecked against the pages by a human.
%
% Pages 70, 72, 74, 76, 78 and 80 are unrelated typed versos and are not
% transcribed. Page 63 carries text on its upper half only. Page 64 is a
% cover sheet in his hand; its words become the section heading.
\documentclass[11pt,a4paper]{article}
\input{../preamble/grothendieck}

\folder{77}
\batch{4}
\pages{61}{80}
\foldertitle{Quadriques affines, extensions quadratiques : notes manuscrites (s.d.).}
\dating{[à partir de 1977]}
\watermark{Édition de démonstration}

\begin{document}

\page{61}
\note{la page reprend au milieu d'un argument commencé avant ce lot : les données $(L, \vartheta, b_0)$, les épaississements $S_0 \subset S_1 \subset S$, la section $\gamma(b_0)$ et les conditions a), b) sont introduites dans les pages qui précèdent.}

D'après b), ces conditions sont équivalentes dans le cas affine. Donc elles sont équivalentes dans le cas général.

\emph{NB} \struck{Pour pouvoir conclure de ces conditions} \add{Supposons que $(L, \vartheta, b_0)$ satisfait à l'une \uncertain{ou l'autre}} \struck{i.e. que la condition} \ill{} l'on ait \struck{\ill{}} \struck{existe $b_1 \in$}
\[
  \vartheta_1 = \gamma(b_0)
\]
quand peut-on conclure qu'il existe $b_1 \in \Gamma(S_1, L_1)$ avec $b_1$ \uncertain{relevant} $b_0$, $b_1^2 = \vartheta_1$ ? \struck{Or les $b_1$ en} \struck{Bien entendu}, les $b_1$ en question sont exactement les $b_1$ qui relèvent $b_0$ (puisque : $\vartheta_1 = \gamma(b_0)$)), \uncertain{et} l'obstruction est dans $H^1(S_0, L_0)$. \uncertain{Prouvons} donc

\emph{Corollaire 1} \uncertain{Pour} que $(L, \vartheta, b_0)$ soit dans l'\uncertain{image essentielle}, il \struck{\ill{}} suffit, et il faut lorsque $H^1(S_0, L_0) = 0$ \struck{\ill{}} (p.~ex. $S_0$ affine, en particulier si $S$ affine) qu'il existe $b_1 \in \Gamma(S_1, L_1)$ tel que
\[
  b_1^2 = \vartheta_1 .
\]

\emph{Corollaire 2} Supposons $S_0$ réduit. Alors le foncteur $\Omega \longmapsto (L, \vartheta)$ est pleinement fidèle, et son image essentielle est formée des couples $(L, \vartheta)$ satisfaisant les \struck{deux} conditions

\note{de « a) » jusqu'au haut de la page 62 (« … par 2 dans $L_1$ ») le passage est biffé de traits obliques ; il est repris au propre plus bas, p. 62--63.}

\struck{a) $\vartheta_0 = \vartheta|S_0$ est de la forme $b_0^2$ ($b_0 \in \Gamma(S_0, L_0)$)} (NB alors $b_0$ est uniquement déterminé, \struck{dans ce cas,} localement sur $S_0$)

\struck{b) $\gamma(b_0) = \vartheta_1$} \struck{ou encore} \struck{\ill{}} \struck{i.e. encore la section} \struck{$\alpha(b_0, \vartheta_1) \in \Gamma(S_0, L_0)$}

\page{62}
\struck{$\alpha(\vartheta_1) \in \Gamma(S_0, L_0)$} \struck{définie par la condition}
\[
  \text{\struck{$\beta(\alpha(\vartheta_1)) = \gamma(b_0) - \vartheta_1$}}
\]
\struck{où}
\[
  \text{\struck{$\beta : L_0 \xrightarrow{\ \sim\ } \operatorname{Ker}(L_1 \to L_0)$}}
\]
\struck{est l'isom. canonique induit par la multiplication par 2 dans $L_1$)}

\struck{\ill{}} que $\vartheta_1 = \vartheta|S_1$ ($\in \Gamma(S_1, L_1^{\otimes 2})$) soit le carré d'une section de $L_1$. (\uncertain{ce qui suffit}, et il faut si $H^1(S_0, L_0) = 0$,) \struck{cette condition} \struck{équivaut} \add{à ce} qu'il existe $b_1 \in \Gamma(S_1, L_1)$ tel que
\[
  b_1^2 = \vartheta_1 .
\]
En particulier, il suffit, et il faut si $H^1(S, L) \to H^1(S_0, L_0)$ est injectif, qu'il existe $b \in \Gamma(S, L)$ et $c \in \Gamma(S, L^{\otimes 2})$ tels que l'on ait
\[
  b^2 - 4c = \vartheta .
\]
\note{le $L$ de $\Gamma(S, L)$ et de $\Gamma(S, L^{\otimes 2})$ est souligné sur la page.}

En effet, la première assertion résulte du fait que si $S_0$ est réduit, la section $\vartheta_0 = \vartheta|S_0$ admet au plus une \emph{\uncertain{racine}}. Le reste suit de façon évidente.

\emph{NB} La condition que $\vartheta_1 = b_1^2$ \uncertain{localement} se \uncertain{décompose} en deux :

a) $\vartheta_0 = b_0^2$ (où $b_0 \in \Gamma(S_0, L_0)$ est uniquement déterminé) i.e. \struck{\ill{}} la classe $\partial(b_0) \in H^1(S_0, \mu_2)$ est nulle !

b) $\vartheta_1 = \gamma(b_0)$, ce qui peut aussi s'écrire sous la forme
\[
  \alpha(\vartheta_1) = 0
\]
où
\[
  \alpha(\vartheta_1) \in \Gamma(S_0, L_0)
\]
est défini par la condition
\[
  \beta(\alpha(\vartheta_1)) = \gamma(b_0) - \vartheta_1
\]
où
\[
  \beta : L_0 \xrightarrow{\ \sim\ } \operatorname{Ker}(L_1 \to L_0)
\]

\page{63}
est l'isom. can. induit par la multiplication par 2 dans $L_1$.

\emph{Corollaire 3} Supposons $S_0$ réduit et normal. Pour que $(L, \vartheta)$ provienne d'un rev. quadratique de $S$, il f. et s. qu'il en soit ainsi aux points \struck{\ill{}} maximaux de $S_0$.

\section{Discriminant des formes quadratiques}
\note{titre pris sur la page 64, feuillet de couverture portant de sa main « discriminant des formes quadratiques » ; la mention « (16p) » au crayon n'est peut-être pas de lui.}

\page{65}
\emph{Discriminant de formes quadratiques.}

Soit $I$ ens. d'indices fini, $A_I = \mathbf{Z}[I \times I]$ l'anneau \struck{\ill{}} des matrices \add{universelles} carrées de type $I$, $AQ_I = \mathbf{Z}[\mathrm{Sym}^2(I)]$ \add{celui des} \struck{\ill{}} matrices carrées symétriques universelles, \struck{\ill{}} dont les générateurs sur $\mathbf{Z}$ sont \struck{$A\mathbf{Z}^{I \times I} \to \mathbf{Z}$} de la forme $Y_J$, $J$ parcourant l'ens. des parties de $I$ de cardinal 1 ou 2. Coord. de $A_I$ \struck{$AQ_I \to AQ_J$} les $X_{ij}$ ($(i,j) \in I \times I$) \struck{\ill{}}. Considérons
\[
  A_I \xrightarrow{\ \varphi\ } AQ_I
\]
défini par
\[
  \varphi(X_{ij}) =
  \begin{cases}
    Y_{\{i,j\}} & \text{si } i \neq j \\
    2Y_{\{i,j\}} = 2X_{\{i\}} & \text{si } i = j
  \end{cases}
\]
\note{les générateurs de $AQ_I$ sont écrits avec une lettre surchargée, $X$ et $Y$ l'un sur l'autre, ici et p. 66 ; on lit $Y$, notation qu'il garde p. 67 et suivantes. Le « $2X_{\{i\}}$ » final est tel quel sur la page.}

qui correspond : à l'homom. de foncteurs
\[
  MS_I(k) \longrightarrow M_I(k)
\]
associant : à une forme quadratique de coeff. les $a_{\{i,j\}}$, la matrice de celle-ci. Soit
\[
  \Delta_0 \in A_I
\]
le déterminant de la $I$-matrice universelle, $\Delta = \varphi(\Delta_0)$ son image, qui est le discriminant de la forme quadratique universelle indexée par $I$.

On a
\[
  \Delta_0 = \sum_{\sigma \in \mathfrak{S}_I} \varepsilon_\sigma
  \underbrace{\prod_{i \in I} X_{i, \sigma i}}_{X_\sigma}
\]
Considérons la décomposition de $I$ en cycles suivant $\sigma$, le produit s'écrit \ill{} en mettant à part

\page{66}
les cycles de longueur 1 et ceux de longueur 2
\[
  X_\sigma = \prod_{i \in I_1(\sigma)} X_{ii}
  \prod_{i \in I_2(\sigma)} X_{\{i, \sigma i\}}
  \prod_{i \in I_*(\sigma)} X_{\{i, \sigma i\}}
\]
\[
  \varepsilon_\sigma = (-1)^{\nu_2 + \nu_*^+}
\]
où
\begin{itemize}
  \item $I_1(\sigma) = I^\sigma$
  \item $I_2(\sigma) = \bigcup$ cycles de card.~2
  \item $I_*(\sigma) = \bigcup$ cycles de card $\geq 3$
  \item $\nu_1 =$ nb de cycles de card 1 $= \operatorname{card} I^\sigma$ \add{$= \operatorname{card} I_1(\sigma)$}
  \item $\nu_2 =$ nb des cycles de card 2
  \item $\nu_*^+ =$ nb des cycles de cardinal pair $\geq 2$
  \item $\nu_* =$ nb des \ill{} $\geq 2$
\end{itemize}
\note{dans les deux dernières lignes le chiffre de la borne est repassé à l'encre forte et se lit 2 ; vu la définition de $I_*(\sigma)$, on attendrait 3.}

On trouve donc
\[
  \Delta = \sum_{\sigma \in \mathfrak{S}_I} (-1)^{\nu_2 + \nu_*^+} 2^{\nu_1}
  \underbrace{\prod_{i \in I_1(\sigma)} Y_{\{i\}}
  \Bigl(\prod_{H \in \dot I_2(\sigma)} Y_H\Bigr)^2
  \prod_{H \in \dot I_*(\sigma)} Y_{H,\sigma}}_{Y_\sigma}
\]
où
\[
  Y_{H,\sigma} = \prod_{a \in A(H,\sigma)} X_a ,
\]
$A(H, \sigma)$ étant l'ens. des \emph{arêtes} du polygone combinatoire (orienté) dont l'ens. des sommets est $H$, avec la structure polyg. orientée définie par $\sigma|H$ ($\dot I_2(\sigma) = I_2(\sigma)/\sigma$, $\dot I_*(\sigma) = I_*(\sigma)/\sigma$).

\struck{Les monômes} \add{On} \uncertain{voit} donc que \add{la donnée de $\sigma$ équivaut à celle d'} \struck{$\sigma$ définit} une décomposition de $I$ en
\[
  I_1(\sigma) \sqcup I_2(\sigma) \sqcup I_3(\sigma),
\]
où $I_2(\sigma)$ est munie d'une involution sans pt fixe $\sigma_2$ (i.e. d'une relation d'équivalence : classes de cardinal 2) \struck{et} et $I_3(\sigma)$ d'une structure de

\page{67}
contour orienté, et
\[
  Y_\sigma = \prod_{i \in I_1(\sigma)} Y_i
  \Bigl(\prod_{a \in I_2(\sigma)/\sigma_2} Y_a\Bigr)^2
  \prod_{a \in A(I_3(\sigma))} Y_a
  \;=\; Y_{\gamma(\sigma)}
\]
\note{sous $A(I_3(\sigma))$ : « arêtes ».}

Quant aux monômes \add{$Y_\sigma$}, ils correspondent 1--1 aux structures de graphes \add{$\gamma$} stricts sur $I$ correspondant à $\sigma$, i.e. dont les composantes connexes sont soit des points isolés, soit des \struck{simplexes} doublets $\bullet\!\!-\!\!\bullet$, soit des polygones combinatoires d'ordre \struck{$d$} $\geq 3$. Pour que \add{dans} $Y_\sigma = Y_{\sigma'}$, il f. et s. que $\sigma$ et $\sigma'$ définissent les \uncertain{mêmes} $I_1$ et $I_2$, \struck{\ill{}} et que la structure \struck{polyg.} de \uncertain{contour} sur $I_3 = I \setminus I_1 \setminus I_2$, dont \uncertain{proviennent} \struck{la partition} \add{\uncertain{les contours}} \struck{partition en} \ill{} de $I_3$ en \struck{\ill{}} \ill{}, sur \uncertain{chacune} \uncertain{desquelles} $\sigma$ et $\sigma'$ sont ou bien \uncertain{égales}, ou bien inverses. Donc chaque $Y_\gamma$ intervient exactement $2^{\nu_*}$ fois. Donc
\[
  \Delta = \sum_{\gamma \in \Gamma_I} (-1)^{\nu_2 + \nu_*^+(\gamma)}
  2^{\nu_1 + \nu_*(\gamma)}
  \underbrace{\prod_{i \in I_1(\gamma)} Y_i
  \Bigl(\prod_{a \in A(I_2(\gamma))} Y_a\Bigr)^2
  \prod_{a \in A(I_3(\gamma))} Y_a}_{M_\gamma}
\]
\note{sous $\gamma \in \Gamma_I$ : « ens. des structures de graphes sur $I$ du type envisagé ».}

\page{68}
\[
  \Delta \bmod 2 = \sum_{\gamma \in \Gamma_0}
  \Bigl(\prod_{a \in A(\gamma_0)} Y_a\Bigr)^2
\]
\note{devant la parenthèse, un symbole biffé. Sous $\gamma \in \Gamma_0$ : « structures de graphes \uncertain{sans} sommets isolés -- i.e. correspondant aux invol. ss pt fixe de $I$ -- i.e. telles que $\nu_1(\gamma) = \nu_*(\gamma) = 0$ ».}

On voit donc que $\Delta \bmod 2 = 0$ si $\operatorname{card} I$ est impair. \struck{Soit $B = \sum$ \ill{}} Supposons pour la suite
\[
  \operatorname{card} I = m = 2n
\]
\struck{Notons} et posons
\[
  B = B_I = \sum_{\gamma \in \Gamma_0} \Bigl(\prod_{a \in A(I, \gamma_2)} Y_a\Bigr) .
\]
\note{entre $\sum$ et la parenthèse, une rature noircie.}

(On a aussi \struck{$n = \nu_2(\gamma) = \ldots$} $\nu_2(\gamma) = n$, $\nu_*^+(\gamma) = 0$, \uncertain{donc} $(-1)^n$ est le coeff avec lequel $\prod Y_a$ figure dans $\Delta$ \ill{}) \struck{\ill{}} Soit $\varepsilon = \pm 1$, \uncertain{question} aussi de déterminer $\varepsilon$ de façon que
\[
  B^2 - \varepsilon \Delta \equiv 0 \quad (4)
\]
i.e. tel que l'on puisse trouver $C \in AQ_I$ avec
\[
  B^2 - 4C = \varepsilon \Delta .
\]
On \struck{voit} \add{prévoit} que $\varepsilon = (-1)^n$.
\[
  B^2 = \sum_{\gamma \in \Gamma_0} \Bigl(\prod_{a \in A(\gamma)} Y_a\Bigr)^2
  + 2 \sum_{\substack{\gamma, \gamma' \in \Gamma_0 \\ \gamma \neq \gamma'}}
  \Bigl(\prod_{a \in A(\gamma)} Y_a\Bigr)\Bigl(\prod_{a' \in A(\gamma')} Y_{a'}\Bigr)
\]
\note{le coefficient du second terme est surchargé (4 et 2 l'un sur l'autre) ; un symbole biffé précède la première parenthèse.}

Donc \ill{}

\page{69}
\note{le haut de la page, jusqu'à « $2^{\nu_1(\gamma) + \nu_*(\gamma)} \equiv 0$ (4) », est biffé d'un long trait horizontal et de traits obliques.}

\struck{$(-1)^n \Delta \equiv B + \sum$, la somme sur $\gamma \in \Gamma_1 =$ ens. des $\gamma \in \Gamma$ tels que $\nu_1(\gamma) + \nu_*(\gamma) = 1$,} \struck{i.e. [$\nu_1(\gamma) = 1$, $\nu_2(\gamma) = 0$ impossible car $m$ \uncertain{pair}] ou $\nu_1(\gamma) = 0$, $\nu_2(\gamma) = 1$ i.e. pas de pt fixe, et \uncertain{seulement} une composante connexe de card $\geq 3$}

\struck{Notons que $\nu_1(\gamma) + \nu_*(\gamma) \equiv m$ (2) puisque $m - (\nu_1(\gamma) + \nu_*(\gamma)) = \nu_2(\gamma) \equiv 0$ (2). Si $m$ pair, on trouve donc $\nu_1(\gamma) + \nu_*(\gamma) \equiv 0$ (2), donc si $\nu_1(\gamma) + \nu_*(\gamma) \neq 0$, on a $\nu_1(\gamma) + \nu_*(\gamma) \geq 2$ donc $2^{\nu_1(\gamma) + \nu_*(\gamma)} \equiv 0$ (4).} \struck{D'où} \ill{}

\[
  \Delta \equiv (-1)^n \sum_{\gamma \in \Gamma_0}
  \Bigl(\prod_{a \in A(\gamma)} Y_a\Bigr)^2
  + 2 \sum_{\gamma \in \Gamma_1} (-1)^{\nu_2(\gamma) + 1}
  \qquad \bmod (4)
\]
\note{au-dessus, un mot biffé ; dans le premier terme, une rature noircie après $(-1)^n$ et une autre devant la parenthèse. Encadré à gauche : « $\Gamma_1$ : graphes $\gamma \in \Gamma$ tels que $\nu_1(\gamma) = 0$, $\nu_*(\gamma) = \nu_*^+(\gamma) = 1$ donc ». La seconde somme s'arrête à l'exposant.}

\[
  B^2 - (-1)^n \Delta \equiv
  \sum_{\substack{\gamma, \gamma' \in \Gamma_0 \\ \gamma \neq \gamma'}}
  \underbrace{\Bigl(\prod_{a \in A(\gamma)} Y_a\Bigr)
  \Bigl(\prod_{a' \in A(\gamma')} Y_{a'}\Bigr)}_{M_{\gamma, \gamma'}}
  \quad \bmod 4
\]

Considérons un couple $\gamma, \gamma'$ fixé, \ill{} \ill{}, 2 \uncertain{invol.} \ill{} sans pts fixes $\sigma, \sigma'$, distinctes. \struck{Ainsi,} Examinons la structure du monôme $M_{\gamma, \gamma'}$ : si nous posons
\[
  I_{\sigma, \sigma'} = \{ i \in I \mid \sigma i = \sigma' i \}, \quad
  I = I_{\sigma, \sigma'} \sqcup J_{\sigma, \sigma'} ,
\]
\uncertain{avec} $J_{\sigma, \sigma'} \neq \varnothing$. Dans

\page{71}
$J_{\sigma, \sigma'}$, \struck{\ill{}} les $a \in A(\gamma)$ et les $a' \in A(\gamma')$ sont distincts. \struck{\ill{}} Le produit des $Y_a$, $Y_{a'}$ correspondants est aussi le monôme correspondant à une structure \struck{polyg. \ill{}} \add{de contour} sur $\mathrm{II}(\sigma, \sigma')$ (dont les arêtes sont les ens. de 2 éléments $\{i, \sigma i\}$ et $\{i, \sigma' i\}$), dont les composantes sont d'ordre pair $\geq 4$. \note{le symbole noté ici $\mathrm{II}$ est une capitale à double hampe ; c'est sans doute le $J_{\sigma,\sigma'}$ de la page 69.}

Le couple $(\sigma, \sigma')$ est déterminé, à symétrie \add{$(\sigma, \sigma') \mapsto (\sigma', \sigma)$} près, \ill{} \uncertain{dès} que \ill{} composantes conn. de $\mathrm{II}(\sigma, \sigma')$, \uncertain{puis} par la structure de graphe correspondante \add{\uncertain{ainsi}} \ill{} -- et on trouve \uncertain{toutes} les structures de graphe $\gamma \in \Gamma$ pour lesquelles \struck{\ill{}} $\nu_1(\gamma) = 0$, et \struck{les cycles} $\nu_*^+(\gamma) = \nu_*(\gamma)$. Un même monôme intervient $2^{\nu_*(\gamma)} = 2^{\nu_*^+(\gamma)}$ fois.

Donc on trouve
\[
  B^2 - (-1)^n \Delta \equiv
  \sum_{\substack{\gamma \in \Gamma \\ \nu_1(\gamma) = 0 \\ \nu_*(\gamma) = \nu_*^+(\gamma) > 0}}
  2^{\nu_*(\gamma)} M_\gamma
  \quad \bmod 4
\]
\[
  + \, 2 \sum_{\gamma \in \Gamma_1} (-1)^{n + \nu_2(\gamma)} M_\gamma
\]
\[
  \equiv 2 \sum_{\gamma \in \Gamma_1} M_\gamma \quad \not\equiv 0 \quad (4)\,!
\]
\[
  + \, 2 \sum (-1)^{\varepsilon(\gamma)} M_\gamma
\]
\[
  = 2 \sum M_\gamma \bigl(\underbrace{1 + (-1)^{\varepsilon(\gamma)}}_{0 \text{ ou } 2}\bigr)
  \equiv 0 \quad (4)
\]
\note{la ligne « $\equiv 2 \sum M_\gamma \not\equiv 0$ (4) ! » est d'une encre brune ; sous son $\sum$, une condition raturée se termine par « $= 1$ ». Les deux lignes suivantes, d'une autre encre, la corrigent.}

\page{73}
\note{nouveau départ, sans titre.}

$A$ anneau

$\Delta \in A$

$A_0 = A / 2A$

On veut écrire

(1)
\[
  \Delta = b^2 - 4c \qquad b, c \in A
\]
ce qui implique
\[
  \Delta_0 \in A_0^2 \qquad (\Delta_0 = (\Delta \bmod 2))
\]
\struck{soit}
\[
  \Delta_0 = b_0^2 \qquad b_0 \in A_0
\]
où $b_0$ est uniquement déterminé si $A_0$ \uncertain{réduit} (a), par la condition $b \equiv b_0$ (2), donc si \ill{} $b$ déterminé mod 2, donc si $b' = b + 2x$, \ill{}. $b'$ satisfait \uncertain{à la} \uncertain{même} condition i.e.
\[
  b'^2 = b^2 + 4(x^2 + bx) \qquad \text{i.e.} \qquad b'^2 \equiv b^2 \quad (4)
\]
\note{« $A_0$ réduit » est encadré et marqué (a) en cercle ; « (4) » est souligné deux fois.}

On sait donc

\emph{Lemme} $\forall\, b_0 \in A_0$, il y a un unique \uncertain{élément} $\gamma(b_0)$ de $A/4A$ tel que $\forall\, b \in A$ \struck{\ill{}} avec $b \equiv b_0$ (2), on ait $\gamma(b_0) \equiv b^2$ (4).

Pour $\Delta \in A$ donné, pour qu'il existe $b, c \in A$ tels que on ait (1), il f. et s. que

a) $\Delta_0$ ($= \Delta \bmod 2$) soit le carré d'un él. $b_0$ de $A_0$ (bien déterminé si $A_0$ réduit)

b) $\gamma(b_0) =$ \uncertain{$\Delta$}.

Supposons qu'on puisse trouver \struck{un} $u \in A^*$ \struck{\ill{}} \struck{$\nu = 2d + 1$}

\struck{tel que les conditions \ill{} soient \ill{}}

\page{75}
\note{sur cette page et les suivantes, les hypothèses marquées (a) à (f) sont encadrées sur la page.}
\note{le haut de la page est très raturé : « \struck{\ill{} tel que} \add{$v \in A$ tel que} $\Delta v^2 =$ \struck{\ill{}} », suivi de plusieurs mots biffés ; puis :}

\add{\uncertain{Ceci implique que}}
\[
  \Delta_0 v_0^2 = b_0'^2
\]
\note{devant $b_0'^2$, plusieurs ratures noircies.}

on aura, si on suppose que $\Delta_0$ est él. régulier \emph{de $A_0$} (b), que
\[
  \Bigl(\frac{b'_0}{v_0}\Bigr)^2 = \Delta_0 \qquad \text{relation dans } A_{0\Delta_0}
\]
\note{au dénominateur, une rature noircie devant $v_0$.}

Supposons que $A_0$ soit intégralement clos (c), alors on aura $b'_0 / v_0 \in A_0$, soit $b_0$, donc
\[
  b'_0 = b_0 v_0 \qquad (b_0 \in A_0,\ b_0^2 = \Delta_0) .
\]
\struck{Ce qui \ill{} \ill{}} Bien entendu, on a alors, si $b \equiv b_0$ (2), $b' = bv$ a \uncertain{même} image : $b'_0$ (2) donc
\[
  \gamma(b'_0) \equiv b'^2 \pmod 4 = b^2 v^2
\]
et il vient : \uncertain{écrivons}
\[
  b^2 v^2 \equiv \Delta v^2 \quad (4)
\]
\note{devant $b^2v^2$ et dans l'exposant de $\Delta$, des ratures noircies ; on devine « $+1$ » dans l'exposant raturé.}

Si en outre, si \struck{$\Delta_0$} $v$ \add{mod 4} est régulier dans $A/4A$ (d), \struck{\ill{}}

2 est régulier dans $A$

$v_0$ est régulier dans $A_0$

\struck{$\Delta_0$ est régulier dans $A_0$}

\struck{on aura donc que}
\[
  b^2 \equiv \Delta \quad (4)
\]
\note{le 2 de l'exposant est surchargé d'un trait.}

\struck{donc}

\page{77}
\marginal{\emph{NB} -- \uncertain{si} $v \in A^*$, \ill{} $v \in A$ tel que \uncertain{$v$ soit régulier} \ill{} et $v_0$ \uncertain{régulier} de $A_0$}

Ainsi on a prouvé

\emph{Prop.} Supposons 2 régulier dans $A$, $A_0 = A/2A$ réduit et intégralement clos, $\Delta \in A$ tel que $\Delta_0 \in A_0$ soit él. régulier de $A_0$. Soit $\nu = 2d + 1$ un \uncertain{entier} \uncertain{naturel} \uncertain{impair}. Pour qu'il \add{existe} $v \in A^*$ \struck{existe} $b', c' \in A$, avec
\[
  \begin{cases}
    b'^2 - 4c' = \Delta^{\nu} v^2 \\
    v_0 \text{ est } A_0\text{-régulier}
  \end{cases}
\]
il f. et s. qu'il existe $b, c \in A$ tels que
\[
  b^2 - 4c = \Delta .
\]
\note{dans l'exposant de $\Delta$, une rature noircie devant $\nu$.}

\struck{Supposons} \add{Considérons} maintenant \struck{l'anneau \ill{}} \add{l'anneau}

(\uncertain{a})
\[
  A_{2\Delta} = (A_2)_\Delta = (A_\Delta)_2 , \quad \text{localisé de } A \text{ pour } 2\Delta
\]
et l'algèbre
\[
  B_{2\Delta} = A_{2\Delta}[T] / (T^2 - \Delta)
\]
Alors on \uncertain{a} \add{\uncertain{les revêtements} \ill{} de groupe $\mathbf{Z}/2\mathbf{Z}$.} C'est une $A_{2\Delta}$-alg. étale \ill{} de $\Delta$. \struck{Supp} \emph{Supposons} que $B_{2\Delta}$ se prolonge \add{\uncertain{finie}} en une algèbre $B_\Delta$ \uncertain{finie} $A_\Delta$ \struck{étale} \struck{\ill{}}

[\uncertain{Exemple} \uncertain{unique} : \uncertain{$\Delta$ nombre premier}, \ill{}, 2 \uncertain{non} régulier dans $A_\Delta$ \ill{}]. Supposons de plus

(e) $\operatorname{Pic} A_\Delta = 0$

\note{au-dessus de l'encadré, un $\mathbf{Z}$ biffé.}

alors on sait (Comme $B_\Delta$ est fini loc. libre de rang 2) que l'on peut écrire

\page{79}
\[
  B_\Delta \simeq A_\Delta[T] / (T^2 + bT + c) \qquad b, c \in A_\Delta
\]
avec l'idéal discriminant de $B_\Delta$ égal à :
\[
  \mathfrak{D}(B_\Delta / A_\Delta) = (b^2 - 4c) A_\Delta
\]
donc, à cause de l'hypothèse $B_\Delta$ étale sur \uncertain{$A_\Delta$},
\[
  b^2 - 4c \in A_\Delta^* .
\]
D'autre part, prenant la restriction à $A_{2\Delta}$, on trouve
\[
  b^2 - 4c = u^2 \Delta , \qquad u \in A_{2\Delta}^* = (A_\Delta)_2^*
\]

\note{de « Supposons que » jusqu'à « $v \in A^*$ » le passage est biffé de traits obliques.}

\struck{Supposons que} $(A_{2\Delta})^* = \{ 2^\mu \Delta^\nu v \mid \mu, \nu \in \mathbf{Z},\ v \in A^* \}$ \struck{(\ill{})}

\struck{dans ce cas}
\[
  \text{\struck{$b^2 - 4c = 2^{2\mu} \Delta^{2\nu + 1} v^2$ dans $(A_{2\Delta})_2$}}
\]
\struck{avec $v \in A^*$}

Comme $b^2 - 4c \in A_\Delta^*$, \ill{} conclut \add{$u^2 \in A_\Delta^*$} \uncertain{ce qui implique} \struck{(\ill{})} $A_\Delta$ intégralement clos dans $A_{2\Delta}$ (f)

$*$ $u \in A_\Delta^*$ \struck{$\mu = 0$}, donc
\[
  b^2 - 4c = \Delta u^2 \qquad u \in A^*
\]
($*$)
\note{dans l'exposant de $\Delta$, une rature noircie.}

Inversement si on a une relation \ill{}, \ill{}

\struck{$\times$ Supposons que $v \in A_\Delta^*$ \ill{}} \struck{les hypothèses \ill{}} \struck{\ill{} d'un \ill{}}
\[
  \text{\struck{$b^2 - 4c = \lambda^2 \Delta \quad \lambda \in A^* \quad (\lambda = (\Delta^\nu v)^2)$}}
\]
\note{l'argument se poursuit au-delà de ce lot.}

\end{document}
